% SPDX-License-Identifier: Apache-2.0
\section{What the format turned out to be}
\label{sec:format}

This section states the structure the corpus established. It is a
summary; the repository's \texttt{docs/format.md} holds the findings
table, roughly 290 rows, each with the command that reproduces it, and
four documents that give the offsets.

\subsection{Container}

A database begins with a fixed header whose first bytes read
\texttt{Xilinx WAVE DATABASE 01} and \texttt{Xilinx Simulator}. Three
pointers at offset \texttt{0x48} locate directory entries, and each
entry is followed by the section it describes: the event section, the
type table, and the debug section. A trailer of \texttt{0x48} bytes
sits before the first directory pointer and holds the simulation end
time, the offset of a marker, the page size and the size of the handle
space.

Values live in \emph{arenas}. An arena spans \texttt{0x800} bytes of
handle space, and the arena table has
$\lceil \mathit{handlespace} / \texttt{0x800} \rceil$ slots. That rule
replaced an earlier reading of the same field, $\lceil n/4 \rceil + 1$
over the object count, which fitted every case until a case with seven
signals broke it. The page directory holds one record of \texttt{0x4c0}
bytes per arena, listing the file offset and compressed length of each
of its pages.

A page is a zlib stream that inflates to exactly 10240 bytes. Inside
it are the time of the first record, the span to the last, a record
count, and then the records, each of which is a time, a key, a length
and the value bytes. There is no alignment between records.

\subsection{Type table}

The type table starts with \texttt{Xilinx ISim TYPE FILE 001} and a
count. Each entry is a length-prefixed name followed by a body whose
shape depends on a kind byte. The kinds observed are enumerations,
SystemVerilog named values, integers with bounds, reals with bounds,
aliases, VHDL access and file types, physical types with their unit
table, arrays with their index constraints, and records with their
fields. Under \texttt{xelab -debug all}, five more appear for the
SystemVerilog string, queue, dynamic array, associative array and
class.

An enumeration entry carries its literals, which is what allows a
decoded value to be rendered as \texttt{'1'}, \texttt{IDLE} or
\texttt{TRUE} rather than as an index, and it is why FST output can
keep the literal.

\subsection{Hierarchy}

The debug section starts with \texttt{Xilinx ISim DBG 006} and
18~region offsets. The regions hold scope records (name, parent,
children, first object, unit, file, line), unit records, declaration
records (name, file, line, size, type, index constraints, kind, port
mode), the source file table, and instance records that bind a
declaration in a scope to a handle.

Two properties of that binding matter to a reader:

\begin{itemize}
  \item Objects connected by a port share one handle. A signal, the
        ports it drives and the ports of the instances below all
        resolve to the same storage, and each contributes one record
        at time~0.
  \item A port bound to a slice of a signal shares the signal's handle
        and carries an offset: in bytes for VHDL, in bits from the low
        end for Verilog. Its value is that window of the signal's
        value.
\end{itemize}

\subsection{Values}

A VHDL value is laid out by type with alignment; a Verilog value is
word pairs, \texttt{[a][b]} per 32 bits, where bit $i$ is
$a_i + 2b_i$ over the four literals \texttt{0}, \texttt{1}, \texttt{Z},
\texttt{X}.

Two rules govern how a value reaches the file.

\emph{Partial writes.} An assignment to a field, an element or a slice
writes only the bytes it touched, keyed at the handle plus the byte
offset of the part. A reader overlays each record on the value the
earlier records built. For Verilog the unit is the word pair.

\emph{Chunking.} A value of 275 bytes or more is split into records.
With $s = \mathit{size} + 24$, the count is
$2 \lceil s/299 \rceil$ and each chunk is
$\lfloor \mathit{size} / \mathit{count} \rfloor$ bytes, the last
taking the remainder, and a remainder of 276 bytes or more is chunked
again by the same rule. The rule was fitted to 55 values from 200 to
30\,000 bytes and holds for every one; the constants 24 and 299 have
no explanation and are recorded as unexplained.

The chunk map belongs to the signal at the handle rather than to the
object being decoded. The distinction is invisible until a port is
bound to a slice of a signal large enough to be chunked, which is what
one of the designs in Section~\ref{sec:validation} produced.
